\Chapter{94}

\Verse{1}
{
    \zR{话说赖大带了贾芹出来，一宿无话，静候贾政回来。}
}
{
    \eR{[English source not present in the checked-in Bencraft/Joly files.]}
}
{
}

\Verse{2}
{
    \zR{单是那些女尼女道重进园 来，都喜欢的了不得，欲要到各处逛逛，明日预备进宫。}
}
{
    \eR{[English source not present in the checked-in Bencraft/Joly files.]}
}
{
}

\Verse{3}
{
    \zR{不料赖大便吩咐了看园的 婆子并小厮看守，惟给了些饭食，却是一步不准走开。}
}
{
    \eR{[English source not present in the checked-in Bencraft/Joly files.]}
}
{
}

\Verse{4}
{
    \zR{那些女孩子摸不着头脑，只 得坐着，等到天亮。}
}
{
    \eR{[English source not present in the checked-in Bencraft/Joly files.]}
}
{
}

\Verse{5}
{
    \zR{园里各处的丫头虽都知道拉进女尼们来，预备宫里使唤，却也 不能深知原委。}
}
{
    \eR{[English source not present in the checked-in Bencraft/Joly files.]}
}
{
}

\Verse{6}
{
    \zR{到了明日早起，贾政正要下班，因堂上发下两省城工估销册子，立刻要查核， 一时不能回家，便叫人回来告诉贾琏，说：“赖大回来，你务必查问明白。}
}
{
    \eR{[English source not present in the checked-in Bencraft/Joly files.]}
}
{
}

\Verse{7}
{
    \zR{该如何 办就如何办了，不必等我。”}
}
{
    \eR{[English source not present in the checked-in Bencraft/Joly files.]}
}
{
}

\Verse{8}
{
    \zR{贾琏奉命，先替芹儿喜欢，又想道：若是办得一点影 儿都没有，又恐贾政生疑，“不如回明二太太，讨个主意办去，便是不合老爷的心， 我也不至甚担干系。”}
}
{
    \eR{[English source not present in the checked-in Bencraft/Joly files.]}
}
{
}

\Verse{9}
{
    \zR{主意定了，进内去见王夫人，陈说：“昨日老爷见了揭帖生 气，把芹儿和女尼女道等都叫进府来查办。}
}
{
    \eR{[English source not present in the checked-in Bencraft/Joly files.]}
}
{
}

\Verse{10}
{
    \zR{今日老爷没空问这件不成体统的事，叫 我来回太太，该怎么便怎么样。}
}
{
    \eR{[English source not present in the checked-in Bencraft/Joly files.]}
}
{
}

\Verse{11}
{
    \zR{我所以来请示太太，这件事如何办理？”}
}
{
    \eR{[English source not present in the checked-in Bencraft/Joly files.]}
}
{
}

\Verse{12}
{
    \zR{王夫人听了诧异道：“这是怎么说!若是芹儿这么样起来，这还成咱们家的人 了么？}
}
{
    \eR{[English source not present in the checked-in Bencraft/Joly files.]}
}
{
}

\Verse{13}
{
    \zR{但只这个贴帖儿的也可恶，这些话可是混嚼说得的么？}
}
{
    \eR{[English source not present in the checked-in Bencraft/Joly files.]}
}
{
}

\Verse{14}
{
    \zR{你到底问了芹儿有这件 事没有呢？”}
}
{
    \eR{[English source not present in the checked-in Bencraft/Joly files.]}
}
{
}

\Verse{15}
{
    \zR{贾琏道：“刚才也问过了。}
}
{
    \eR{[English source not present in the checked-in Bencraft/Joly files.]}
}
{
}

\Verse{16}
{
    \zR{太太想，别说他干了没有，就是干了，一 个人干了混账事也肯应承么？}
}
{
    \eR{[English source not present in the checked-in Bencraft/Joly files.]}
}
{
}

\Verse{17}
{
    \zR{但只我想芹儿也不敢行此事：知道那些女孩子都是娘 娘一时要叫的，倘或闹出事来，怎么样呢？}
}
{
    \eR{[English source not present in the checked-in Bencraft/Joly files.]}
}
{
}

\Verse{18}
{
    \zR{依侄儿的主见，要问也不难，若问出来， 太太怎么个办法呢？”}
}
{
    \eR{[English source not present in the checked-in Bencraft/Joly files.]}
}
{
}

\Verse{19}
{
    \zR{王夫人道：“如今那些女孩子在那里？”}
}
{
    \eR{[English source not present in the checked-in Bencraft/Joly files.]}
}
{
}

\Verse{20}
{
    \zR{贾琏道：“都在园 里锁着呢。”}
}
{
    \eR{[English source not present in the checked-in Bencraft/Joly files.]}
}
{
}

\Verse{21}
{
    \zR{王夫人道：“姑娘们知道不知道？”}
}
{
    \eR{[English source not present in the checked-in Bencraft/Joly files.]}
}
{
}

\Verse{22}
{
    \zR{贾琏道：“大约姑娘们也都知道 是预备宫里头的话，外头并没提起别的来。”}
}
{
    \eR{[English source not present in the checked-in Bencraft/Joly files.]}
}
{
}

\Verse{23}
{
    \zR{王夫人道：“很是。}
}
{
    \eR{[English source not present in the checked-in Bencraft/Joly files.]}
}
{
}

\Verse{24}
{
    \zR{这些东西一刻也 是留不得的。}
}
{
    \eR{[English source not present in the checked-in Bencraft/Joly files.]}
}
{
}

\Verse{25}
{
    \zR{头里我原要打发他们去来着，都是你们说留着好，如今不是弄出事来 了么？}
}
{
    \eR{[English source not present in the checked-in Bencraft/Joly files.]}
}
{
}

\Verse{26}
{
    \zR{你竟叫赖大带了去细细儿的问他的本家儿有人没有，将文书查出，花上几十 两银子，雇只船，派个妥当人，送到本地，一概连文书发还了，也落得无事。}
}
{
    \eR{[English source not present in the checked-in Bencraft/Joly files.]}
}
{
}

\Verse{27}
{
    \zR{若是 为着一两个不好，个个都押着他们还俗，那又太造孽了。}
}
{
    \eR{[English source not present in the checked-in Bencraft/Joly files.]}
}
{
}

\Verse{28}
{
    \zR{若在这里发给官媒，虽然 我们不要身价，他们弄去卖钱，那里顾人的死活呢？}
}
{
    \eR{[English source not present in the checked-in Bencraft/Joly files.]}
}
{
}

\Verse{29}
{
    \zR{芹儿呢，你便狠狠的说他一顿， 除了祭祀喜庆，无事叫他不用到这里来。}
}
{
    \eR{[English source not present in the checked-in Bencraft/Joly files.]}
}
{
}

\Verse{30}
{
    \zR{看仔细碰在老爷气头儿上，那可就吃不了 兜着走了。}
}
{
    \eR{[English source not present in the checked-in Bencraft/Joly files.]}
}
{
}

\Verse{31}
{
    \zR{也说给账房儿里，把这一项钱粮档子销了。}
}
{
    \eR{[English source not present in the checked-in Bencraft/Joly files.]}
}
{
}

\Verse{32}
{
    \zR{还打发个人到水月庵，说老 爷的谕，除了上坟烧纸，要有本家爷们到他那里去，不许接待。}
}
{
    \eR{[English source not present in the checked-in Bencraft/Joly files.]}
}
{
}

\Verse{33}
{
    \zR{若再有一点不好风 声，连老姑子一块儿撵出去。”}
}
{
    \eR{[English source not present in the checked-in Bencraft/Joly files.]}
}
{
}

\Verse{34}
{
    \zR{贾琏一一答应了。}
}
{
    \eR{[English source not present in the checked-in Bencraft/Joly files.]}
}
{
}

\Verse{35}
{
    \zR{出去将王夫人的话告诉赖大，说：“太太的主意，叫你这么 办。}
}
{
    \eR{[English source not present in the checked-in Bencraft/Joly files.]}
}
{
}

\Verse{36}
{
    \zR{办完了，告诉我去回太太。}
}
{
    \eR{[English source not present in the checked-in Bencraft/Joly files.]}
}
{
}

\Verse{37}
{
    \zR{你快办去罢。}
}
{
    \eR{[English source not present in the checked-in Bencraft/Joly files.]}
}
{
}

\Verse{38}
{
    \zR{回来老爷来，你也按着太太的话回去。”}
}
{
    \eR{[English source not present in the checked-in Bencraft/Joly files.]}
}
{
}

\Verse{39}
{
    \zR{赖大听说，便道：“我们太太真正是个佛心。}
}
{
    \eR{[English source not present in the checked-in Bencraft/Joly files.]}
}
{
}

\Verse{40}
{
    \zR{这班东西还着人送回去，既是太太好 心，不得不挑个好人。}
}
{
    \eR{[English source not present in the checked-in Bencraft/Joly files.]}
}
{
}

\Verse{41}
{
    \zR{芹哥儿竟交给二爷开发了罢。}
}
{
    \eR{[English source not present in the checked-in Bencraft/Joly files.]}
}
{
}

\Verse{42}
{
    \zR{那贴帖儿的，奴才想法儿查出 来，重重的收拾他才好。”}
}
{
    \eR{[English source not present in the checked-in Bencraft/Joly files.]}
}
{
}

\Verse{43}
{
    \zR{贾琏点头说：“是了。”}
}
{
    \eR{[English source not present in the checked-in Bencraft/Joly files.]}
}
{
}

\Verse{44}
{
    \zR{即刻将贾芹发落。}
}
{
    \eR{[English source not present in the checked-in Bencraft/Joly files.]}
}
{
}

\Verse{45}
{
    \zR{赖大也赶着 把女尼等领出，按着主意办去了。}
}
{
    \eR{[English source not present in the checked-in Bencraft/Joly files.]}
}
{
}

\Verse{46}
{
    \zR{晚上贾政回来，贾琏赖大回明贾政，贾政本是省 事的人，听了也便撂开手了。}
}
{
    \eR{[English source not present in the checked-in Bencraft/Joly files.]}
}
{
}

\Verse{47}
{
    \zR{独有那些无赖之徒，听得贾府发出二十四个女孩子来， 那个不想？}
}
{
    \eR{[English source not present in the checked-in Bencraft/Joly files.]}
}
{
}

\Verse{48}
{
    \zR{究竟那些人能够回家不能，未知着落，亦难虚拟。}
}
{
    \eR{[English source not present in the checked-in Bencraft/Joly files.]}
}
{
}

\Verse{49}
{
    \zR{且说紫鹃因黛玉渐好，园中无事，听见女尼等预备宫内使唤，不知何事便到贾 母那边打听打听。}
}
{
    \eR{[English source not present in the checked-in Bencraft/Joly files.]}
}
{
}

\Verse{50}
{
    \zR{恰遇着鸳鸯下来闲着，坐下说闲话儿，提起女尼的事，鸳鸯诧异 道：“我并没有听见。}
}
{
    \eR{[English source not present in the checked-in Bencraft/Joly files.]}
}
{
}

\Verse{51}
{
    \zR{回来问问二奶奶就知道了。”}
}
{
    \eR{[English source not present in the checked-in Bencraft/Joly files.]}
}
{
}

\Verse{52}
{
    \zR{正说着，只见傅试家两个女人 过来请贾母的安，鸳鸯要陪了上去。}
}
{
    \eR{[English source not present in the checked-in Bencraft/Joly files.]}
}
{
}

\Verse{53}
{
    \zR{那两个女人因贾母正睡晌觉，就与鸳鸯说了一 声儿，回去了。}
}
{
    \eR{[English source not present in the checked-in Bencraft/Joly files.]}
}
{
}

\Verse{54}
{
    \zR{紫鹃问：“这是谁家差来的？”}
}
{
    \eR{[English source not present in the checked-in Bencraft/Joly files.]}
}
{
}

\Verse{55}
{
    \zR{鸳鸯道：“好讨人嫌!家里有了一 个女孩儿，长的好些儿，就献宝的似的，常在老太太跟前夸他们姑娘怎么长的好，
        心地儿怎么好，‘礼貌上又好，说话儿又简绝，做活计儿手儿又巧，会写会算，尊 长上头最孝敬的，就是待下人也是极和平的。’}
}
{
    \eR{[English source not present in the checked-in Bencraft/Joly files.]}
}
{
}

\Verse{56}
{
    \zR{来了就编这么一大套，常说给老太 太听。}
}
{
    \eR{[English source not present in the checked-in Bencraft/Joly files.]}
}
{
}

\Verse{57}
{
    \zR{我听着很烦。}
}
{
    \eR{[English source not present in the checked-in Bencraft/Joly files.]}
}
{
}

\Verse{58}
{
    \zR{这几个老婆子真讨人嫌，我们老太太偏爱听那些个话。}
}
{
    \eR{[English source not present in the checked-in Bencraft/Joly files.]}
}
{
}

\Verse{59}
{
    \zR{老太太 也罢了，还有宝玉，素常见了老婆子便很厌烦的，偏见了他们家的老婆子就不厌烦， 你说奇不奇？}
}
{
    \eR{[English source not present in the checked-in Bencraft/Joly files.]}
}
{
}

\Verse{60}
{
    \zR{前儿还来说：他们姑娘现有多少人家儿来求亲，他们老爷总不肯应， 心里只要和咱们这样人家作亲才肯。}
}
{
    \eR{[English source not present in the checked-in Bencraft/Joly files.]}
}
{
}

\Verse{61}
{
    \zR{夸奖一回，奉承一回，把老太太的心都说活 了。”}
}
{
    \eR{[English source not present in the checked-in Bencraft/Joly files.]}
}
{
}

\Verse{62}
{
    \zR{紫鹃听了一呆，便假意道：“若老太太喜欢，为什么不就给宝玉定了呢？”}
}
{
    \eR{[English source not present in the checked-in Bencraft/Joly files.]}
}
{
}

\Verse{63}
{
    \zR{鸳 鸯正要说出原故，听见上头说：“老太太醒了。”}
}
{
    \eR{[English source not present in the checked-in Bencraft/Joly files.]}
}
{
}

\Verse{64}
{
    \zR{鸳鸯赶着上去，紫鹃只得起身出 来。}
}
{
    \eR{[English source not present in the checked-in Bencraft/Joly files.]}
}
{
}

\Verse{65}
{
    \zR{回到园里，一头走，一头想道：“天下莫非只有一个宝玉？}
}
{
    \eR{[English source not present in the checked-in Bencraft/Joly files.]}
}
{
}

\Verse{66}
{
    \zR{你也想他，我也想 他。}
}
{
    \eR{[English source not present in the checked-in Bencraft/Joly files.]}
}
{
}

\Verse{67}
{
    \zR{我们家的那一位，越发痴心起来了!看他的那个神情儿，是一定在宝玉身上的 了，三番两次的病，可不是为着这个是什么？}
}
{
    \eR{[English source not present in the checked-in Bencraft/Joly files.]}
}
{
}

\Verse{68}
{
    \zR{这家里‘金’的‘银’的还闹不清， 再添上一个什么傅姑娘，更了不得了。}
}
{
    \eR{[English source not present in the checked-in Bencraft/Joly files.]}
}
{
}

\Verse{69}
{
    \zR{我看宝玉的心也在我们那一位的身上啊，听 着鸳鸯的话，竟是见一个爱一个的。}
}
{
    \eR{[English source not present in the checked-in Bencraft/Joly files.]}
}
{
}

\Verse{70}
{
    \zR{这不是我们姑娘白操了心了吗？”}
}
{
    \eR{[English source not present in the checked-in Bencraft/Joly files.]}
}
{
}

\Verse{71}
{
    \zR{紫鹃本是想 着黛玉，往下一想，连自己也不得主意了，不免神都痴了。}
}
{
    \eR{[English source not present in the checked-in Bencraft/Joly files.]}
}
{
}

\Verse{72}
{
    \zR{要想叫黛玉不用瞎操心 呢，又恐怕他烦恼；要是看着他这样，又可怜见儿的。}
}
{
    \eR{[English source not present in the checked-in Bencraft/Joly files.]}
}
{
}

\Verse{73}
{
    \zR{左思右想，一时烦躁起来， 自己啐自己道：“你替人耽什么忧!就是林姑娘真配了宝玉，他的那性情儿也是难 伏侍的。}
}
{
    \eR{[English source not present in the checked-in Bencraft/Joly files.]}
}
{
}

\Verse{74}
{
    \zR{宝玉性情虽好，又是贪多嚼不烂的。}
}
{
    \eR{[English source not present in the checked-in Bencraft/Joly files.]}
}
{
}

\Verse{75}
{
    \zR{我倒劝人不必瞎操心，我自己才是瞎 操心呢，从今以后，我尽我的心伏侍姑娘，其馀的事全不管。”}
}
{
    \eR{[English source not present in the checked-in Bencraft/Joly files.]}
}
{
}

\Verse{76}
{
    \zR{这么一想，心里倒 觉清净。}
}
{
    \eR{[English source not present in the checked-in Bencraft/Joly files.]}
}
{
}

\Verse{77}
{
    \zR{回到潇湘馆来，见黛玉独自一人坐在炕上，理从前做过的诗文词稿。}
}
{
    \eR{[English source not present in the checked-in Bencraft/Joly files.]}
}
{
}

\Verse{78}
{
    \zR{抬头 见紫鹃进来，便问：“你到那里去了？”}
}
{
    \eR{[English source not present in the checked-in Bencraft/Joly files.]}
}
{
}

\Verse{79}
{
    \zR{紫鹃道：“今儿瞧了瞧姐妹们去。”}
}
{
    \eR{[English source not present in the checked-in Bencraft/Joly files.]}
}
{
}

\Verse{80}
{
    \zR{黛玉 道：“可是找袭人姐姐去么？”}
}
{
    \eR{[English source not present in the checked-in Bencraft/Joly files.]}
}
{
}

\Verse{81}
{
    \zR{紫鹃道：“我找他做什么？”}
}
{
    \eR{[English source not present in the checked-in Bencraft/Joly files.]}
}
{
}

\Verse{82}
{
    \zR{黛玉一想：“这话怎 么顺嘴说出来了呢？”}
}
{
    \eR{[English source not present in the checked-in Bencraft/Joly files.]}
}
{
}

\Verse{83}
{
    \zR{反觉不好意思，便啐道：“你找不找与我什么相干!倒茶去 罢。”}
}
{
    \eR{[English source not present in the checked-in Bencraft/Joly files.]}
}
{
}

\Verse{84}
{
    \zR{紫鹃也心里暗笑，出来倒茶。}
}
{
    \eR{[English source not present in the checked-in Bencraft/Joly files.]}
}
{
}

\Verse{85}
{
    \zR{只听园里一叠声乱嚷，不知何故。}
}
{
    \eR{[English source not present in the checked-in Bencraft/Joly files.]}
}
{
}

\Verse{86}
{
    \zR{一面倒茶，一 面叫人去打听。}
}
{
    \eR{[English source not present in the checked-in Bencraft/Joly files.]}
}
{
}

\Verse{87}
{
    \zR{回来说道：“怡红院里的海棠本来萎了几棵，也没人去浇灌他。}
}
{
    \eR{[English source not present in the checked-in Bencraft/Joly files.]}
}
{
}

\Verse{88}
{
    \zR{昨 日宝玉走去瞧，见枝头上好像有了？}
}
{
    \eR{[English source not present in the checked-in Bencraft/Joly files.]}
}
{
}

\Verse{89}
{
    \zR{朵儿似的。}
}
{
    \eR{[English source not present in the checked-in Bencraft/Joly files.]}
}
{
}

\Verse{90}
{
    \zR{人都不信，没有理他。}
}
{
    \eR{[English source not present in the checked-in Bencraft/Joly files.]}
}
{
}

\Verse{91}
{
    \zR{忽然今日开 的很好的海棠花，众人诧异，都争着去看，连老太太、太太都哄动了，来瞧花儿呢。}
}
{
    \eR{[English source not present in the checked-in Bencraft/Joly files.]}
}
{
}

\Verse{92}
{
    \zR{所以大奶奶叫人收拾园里的树叶子，这些人在那里传唤。”}
}
{
    \eR{[English source not present in the checked-in Bencraft/Joly files.]}
}
{
}

\Verse{93}
{
    \zR{黛玉也听见了，知道老 太太来，便更了衣，叫雪雁去打听：“若是老太太来了，即来告诉我。”}
}
{
    \eR{[English source not present in the checked-in Bencraft/Joly files.]}
}
{
}

\Verse{94}
{
    \zR{雪雁去不 多时，便跑来说：“老太太、太太好些人都来了，请姑娘就去罢。”}
}
{
    \eR{[English source not present in the checked-in Bencraft/Joly files.]}
}
{
}

\Verse{95}
{
    \zR{黛玉略自照了 一照镜子，掠了一掠鬓发，便扶着紫鹃到怡红院来，已见老太太坐在宝玉常卧的榻 上。}
}
{
    \eR{[English source not present in the checked-in Bencraft/Joly files.]}
}
{
}

\Verse{96}
{
    \zR{黛玉便说道：“请老太太安。”}
}
{
    \eR{[English source not present in the checked-in Bencraft/Joly files.]}
}
{
}

\Verse{97}
{
    \zR{退后便见了邢王二夫人，回来与李纨、探春、 惜春、邢岫烟彼此问了好。}
}
{
    \eR{[English source not present in the checked-in Bencraft/Joly files.]}
}
{
}

\Verse{98}
{
    \zR{只有凤姐因病未来；史湘云因他叔叔调任回京，接了家 去；薛宝琴跟他姐姐家去住了；李家姐妹因见园内多事，李婶娘带了在外居住：所 以黛玉今日见的只有数人。}
}
{
    \eR{[English source not present in the checked-in Bencraft/Joly files.]}
}
{
}

\Verse{99}
{
    \zR{大家说笑了一回，讲究这花开得古怪。}
}
{
    \eR{[English source not present in the checked-in Bencraft/Joly files.]}
}
{
}

\Verse{100}
{
    \zR{贾母道：“这花儿应在三月里开的，如 今虽是十一月，因节气迟，还算十月，应着小阳春的天气，因为和暖，开花也是有 的。”}
}
{
    \eR{[English source not present in the checked-in Bencraft/Joly files.]}
}
{
}

\Verse{101}
{
    \zR{王夫人道：“老太太见的多，说得是，也不为奇。”}
}
{
    \eR{[English source not present in the checked-in Bencraft/Joly files.]}
}
{
}

\Verse{102}
{
    \zR{邢夫人道：“我听见这 花已经萎了一年，怎么这回不应时候儿开了？}
}
{
    \eR{[English source not present in the checked-in Bencraft/Joly files.]}
}
{
}

\Verse{103}
{
    \zR{必有个原故。”}
}
{
    \eR{[English source not present in the checked-in Bencraft/Joly files.]}
}
{
}

\Verse{104}
{
    \zR{李纨笑道：“老太太 和太太说的都是。}
}
{
    \eR{[English source not present in the checked-in Bencraft/Joly files.]}
}
{
}

\Verse{105}
{
    \zR{据我的糊涂想头，必是宝玉有喜事来了，此花先来报信。”}
}
{
    \eR{[English source not present in the checked-in Bencraft/Joly files.]}
}
{
}

\Verse{106}
{
    \zR{探春 虽不言语，心里想道：“必非好兆。}
}
{
    \eR{[English source not present in the checked-in Bencraft/Joly files.]}
}
{
}

\Verse{107}
{
    \zR{大凡顺者昌，逆者亡；草木知运，不时而发， 必是妖孽。”}
}
{
    \eR{[English source not present in the checked-in Bencraft/Joly files.]}
}
{
}

\Verse{108}
{
    \zR{但只不好说出来。}
}
{
    \eR{[English source not present in the checked-in Bencraft/Joly files.]}
}
{
}

\Verse{109}
{
    \zR{独有黛玉听说是喜事，心里触动，便高兴说道：“当 初田家有荆树一棵，弟兄三个因分了家，那荆树便枯了。}
}
{
    \eR{[English source not present in the checked-in Bencraft/Joly files.]}
}
{
}

\Verse{110}
{
    \zR{后来感动了他弟兄们，仍 旧归在一处，那荆树也就荣了。}
}
{
    \eR{[English source not present in the checked-in Bencraft/Joly files.]}
}
{
}

\Verse{111}
{
    \zR{可知草木也随人的。}
}
{
    \eR{[English source not present in the checked-in Bencraft/Joly files.]}
}
{
}
\ChapterImage{img/chapters/189第94回 宴海棠賈母賞花妖.jpg}


\Verse{112}
{
    \zR{如今二哥哥认真念书，舅舅喜 欢，那棵树也就发了。”}
}
{
    \eR{[English source not present in the checked-in Bencraft/Joly files.]}
}
{
}

\Verse{113}
{
    \zR{贾母王夫人听了喜欢，便说：“林姑娘比方得有理，很有 意思。”}
}
{
    \eR{[English source not present in the checked-in Bencraft/Joly files.]}
}
{
}

\Verse{114}
{
    \zR{正说着，贾赦、贾政、贾环、贾兰都进来看花。}
}
{
    \eR{[English source not present in the checked-in Bencraft/Joly files.]}
}
{
}

\Verse{115}
{
    \zR{贾赦便说：“据我的主意，把 他砍去。}
}
{
    \eR{[English source not present in the checked-in Bencraft/Joly files.]}
}
{
}

\Verse{116}
{
    \zR{必是花妖作怪。”}
}
{
    \eR{[English source not present in the checked-in Bencraft/Joly files.]}
}
{
}

\Verse{117}
{
    \zR{贾政道：“见怪不怪，其怪自败。}
}
{
    \eR{[English source not present in the checked-in Bencraft/Joly files.]}
}
{
}

\Verse{118}
{
    \zR{不用砍他，随他去就 是了。”}
}
{
    \eR{[English source not present in the checked-in Bencraft/Joly files.]}
}
{
}

\Verse{119}
{
    \zR{贾母听见，便说：“谁在这里混说？}
}
{
    \eR{[English source not present in the checked-in Bencraft/Joly files.]}
}
{
}

\Verse{120}
{
    \zR{人家有喜事好处，什么怪不怪的。}
}
{
    \eR{[English source not present in the checked-in Bencraft/Joly files.]}
}
{
}

\Verse{121}
{
    \zR{若 有好事，你们享去；若是不好，我一个人当去。}
}
{
    \eR{[English source not present in the checked-in Bencraft/Joly files.]}
}
{
}

\Verse{122}
{
    \zR{你们不许混说！”}
}
{
    \eR{[English source not present in the checked-in Bencraft/Joly files.]}
}
{
}

\Verse{123}
{
    \zR{贾政听了，不敢 言语，讪讪的同贾赦等走了出来。}
}
{
    \eR{[English source not present in the checked-in Bencraft/Joly files.]}
}
{
}

\Verse{124}
{
    \zR{那贾母高兴，叫人传话到厨房里：“快快预备酒席，大家赏花。”}
}
{
    \eR{[English source not present in the checked-in Bencraft/Joly files.]}
}
{
}

\Verse{125}
{
    \zR{叫宝玉、环 儿、兰儿：“各人做一首诗志喜。}
}
{
    \eR{[English source not present in the checked-in Bencraft/Joly files.]}
}
{
}

\Verse{126}
{
    \zR{林姑娘的病才好，别叫他费心，若高兴，给你们 改改。”}
}
{
    \eR{[English source not present in the checked-in Bencraft/Joly files.]}
}
{
}

\Verse{127}
{
    \zR{对着李纨道：“你们都陪我喝酒。”}
}
{
    \eR{[English source not present in the checked-in Bencraft/Joly files.]}
}
{
}

\Verse{128}
{
    \zR{李纨答应了是，便笑对探春笑道：“都 是你闹的。”}
}
{
    \eR{[English source not present in the checked-in Bencraft/Joly files.]}
}
{
}

\Verse{129}
{
    \zR{探春道：“饶不叫我们做诗，怎么我们闹的？”}
}
{
    \eR{[English source not present in the checked-in Bencraft/Joly files.]}
}
{
}

\Verse{130}
{
    \zR{李纨道：“海棠社不 是你起的么？}
}
{
    \eR{[English source not present in the checked-in Bencraft/Joly files.]}
}
{
}

\Verse{131}
{
    \zR{如今那棵海棠也要来入社了。”}
}
{
    \eR{[English source not present in the checked-in Bencraft/Joly files.]}
}
{
}

\Verse{132}
{
    \zR{大家听着都笑了。}
}
{
    \eR{[English source not present in the checked-in Bencraft/Joly files.]}
}
{
}

\Verse{133}
{
    \zR{一时摆上酒菜，一面喝着，彼此都要讨老太太的喜欢，大家说些兴头话。}
}
{
    \eR{[English source not present in the checked-in Bencraft/Joly files.]}
}
{
}

\Verse{134}
{
    \zR{宝玉 上来斟了酒，便立成了四句诗，写出来念与贾母听，道： 海棠何事忽摧？？}
}
{
    \eR{[English source not present in the checked-in Bencraft/Joly files.]}
}
{
}

\Verse{135}
{
    \zR{今日繁花为底开？}
}
{
    \eR{[English source not present in the checked-in Bencraft/Joly files.]}
}
{
}

\Verse{136}
{
    \zR{应是北堂增寿考，一阳旋复占先梅。}
}
{
    \eR{[English source not present in the checked-in Bencraft/Joly files.]}
}
{
}

\Verse{137}
{
    \zR{贾环也写了来，念道： 草木逢春当茁芽，海棠未发候偏差。}
}
{
    \eR{[English source not present in the checked-in Bencraft/Joly files.]}
}
{
}

\Verse{138}
{
    \zR{人间奇事知多少，冬月开花独我家。}
}
{
    \eR{[English source not present in the checked-in Bencraft/Joly files.]}
}
{
}

\Verse{139}
{
    \zR{贾兰恭楷誊正，呈与贾母。}
}
{
    \eR{[English source not present in the checked-in Bencraft/Joly files.]}
}
{
}

\Verse{140}
{
    \zR{贾母命李纨念道： 烟凝媚色春前萎，霜？}
}
{
    \eR{[English source not present in the checked-in Bencraft/Joly files.]}
}
{
}

\Verse{141}
{
    \zR{微红雪后开。}
}
{
    \eR{[English source not present in the checked-in Bencraft/Joly files.]}
}
{
}

\Verse{142}
{
    \zR{莫道此花知识浅，欣荣预佐合欢杯。}
}
{
    \eR{[English source not present in the checked-in Bencraft/Joly files.]}
}
{
}

\Verse{143}
{
    \zR{贾母听毕，便说：“我不大懂诗，听去倒是兰儿的好，环儿做的不好。}
}
{
    \eR{[English source not present in the checked-in Bencraft/Joly files.]}
}
{
}

\Verse{144}
{
    \zR{都上来吃饭 罢。”}
}
{
    \eR{[English source not present in the checked-in Bencraft/Joly files.]}
}
{
}

\Verse{145}
{
    \zR{宝玉看见贾母喜欢，更是兴头，因想起：“晴雯死的那年，海棠死的；今日 海棠复荣，我们院内这些人，自然都好，但是晴雯不能像花的死而复生了。”}
}
{
    \eR{[English source not present in the checked-in Bencraft/Joly files.]}
}
{
}

\Verse{146}
{
    \zR{顿觉 转喜为悲。}
}
{
    \eR{[English source not present in the checked-in Bencraft/Joly files.]}
}
{
}

\Verse{147}
{
    \zR{忽又想起前日巧姐提凤姐要把五儿补入，“或此花为他而开，也未可知。”}
}
{
    \eR{[English source not present in the checked-in Bencraft/Joly files.]}
}
{
}

\Verse{148}
{
    \zR{却又转悲为喜，依旧说笑。}
}
{
    \eR{[English source not present in the checked-in Bencraft/Joly files.]}
}
{
}

\Verse{149}
{
    \zR{贾母还坐了半天，然后扶了珍珠回去了，王夫人等跟着过来。}
}
{
    \eR{[English source not present in the checked-in Bencraft/Joly files.]}
}
{
}

\Verse{150}
{
    \zR{只见平儿笑嘻嘻 的迎上来，说：“我们奶奶知道老太太在这里赏花，自己不得来，叫奴才来伏侍老 太太、太太们。}
}
{
    \eR{[English source not present in the checked-in Bencraft/Joly files.]}
}
{
}

\Verse{151}
{
    \zR{还有两匹红送给宝二爷包裹这花，当作贺礼。”}
}
{
    \eR{[English source not present in the checked-in Bencraft/Joly files.]}
}
{
}

\Verse{152}
{
    \zR{袭人过来接了，呈 与贾母看。}
}
{
    \eR{[English source not present in the checked-in Bencraft/Joly files.]}
}
{
}

\Verse{153}
{
    \zR{贾母笑道：“偏是凤丫头行出点事儿来，叫人看着又体面，又新鲜，很 有趣儿。”}
}
{
    \eR{[English source not present in the checked-in Bencraft/Joly files.]}
}
{
}

\Verse{154}
{
    \zR{袭人笑着向平儿道：“回去替宝二爷给二奶奶道谢：要有喜，大家喜。”}
}
{
    \eR{[English source not present in the checked-in Bencraft/Joly files.]}
}
{
}

\Verse{155}
{
    \zR{贾母听了，笑道：“嗳哟!我还忘了呢。}
}
{
    \eR{[English source not present in the checked-in Bencraft/Joly files.]}
}
{
}

\Verse{156}
{
    \zR{凤丫头虽病着，还是他想的到，送的也巧。”}
}
{
    \eR{[English source not present in the checked-in Bencraft/Joly files.]}
}
{
}

\Verse{157}
{
    \zR{一面说着，众人就随着去了。}
}
{
    \eR{[English source not present in the checked-in Bencraft/Joly files.]}
}
{
}

\Verse{158}
{
    \zR{平儿私与袭人道：“奶奶说，这花儿开的怪，叫你铰 块红绸子挂挂，就应在喜事上去了。}
}
{
    \eR{[English source not present in the checked-in Bencraft/Joly files.]}
}
{
}

\Verse{159}
{
    \zR{以后也不必只管当作奇事混说。”}
}
{
    \eR{[English source not present in the checked-in Bencraft/Joly files.]}
}
{
}

\Verse{160}
{
    \zR{袭人点头答 应，送了平儿出去不提。}
}
{
    \eR{[English source not present in the checked-in Bencraft/Joly files.]}
}
{
}

\Verse{161}
{
    \zR{且说那日宝玉本来穿着一裹圆的皮袄在家歇息，因见花开，只管出来看一回、 赏一回、叹一回、爱一回的，心中无数悲喜离合，都弄到这株花上去了。}
}
{
    \eR{[English source not present in the checked-in Bencraft/Joly files.]}
}
{
}

\Verse{162}
{
    \zR{忽然听说 贾母要来，便去换了一件狐腋箭袖，罩一件玄狐腿外褂，出来迎接贾母。}
}
{
    \eR{[English source not present in the checked-in Bencraft/Joly files.]}
}
{
}

\Verse{163}
{
    \zR{匆匆穿换， 未将“通灵宝玉”挂上。}
}
{
    \eR{[English source not present in the checked-in Bencraft/Joly files.]}
}
{
}

\Verse{164}
{
    \zR{及至后来贾母去了，仍旧换衣，袭人见宝玉脖子上没有挂 着，便问：“那块玉呢？”}
}
{
    \eR{[English source not present in the checked-in Bencraft/Joly files.]}
}
{
}

\Verse{165}
{
    \zR{宝玉道：“刚才忙乱换衣，摘下来放在炕桌上，我没有 带。”}
}
{
    \eR{[English source not present in the checked-in Bencraft/Joly files.]}
}
{
}

\Verse{166}
{
    \zR{袭人回看桌上，并没有玉，便向各处找寻，踪影全无，吓得袭人满身冷汗。}
}
{
    \eR{[English source not present in the checked-in Bencraft/Joly files.]}
}
{
}

\Verse{167}
{
    \zR{宝玉道：“不用着急，少不得在屋里的。}
}
{
    \eR{[English source not present in the checked-in Bencraft/Joly files.]}
}
{
}

\Verse{168}
{
    \zR{问他们就知道了。”}
}
{
    \eR{[English source not present in the checked-in Bencraft/Joly files.]}
}
{
}

\Verse{169}
{
    \zR{袭人当作麝月等藏起 吓他玩，便向麝月等笑着说道：“小蹄子们，玩呢，到底有个玩法。}
}
{
    \eR{[English source not present in the checked-in Bencraft/Joly files.]}
}
{
}

\Verse{170}
{
    \zR{把这件东西藏 在那里了？}
}
{
    \eR{[English source not present in the checked-in Bencraft/Joly files.]}
}
{
}

\Verse{171}
{
    \zR{别真弄丢了，那可就大家活不成了！”}
}
{
    \eR{[English source not present in the checked-in Bencraft/Joly files.]}
}
{
}

\Verse{172}
{
    \zR{麝月等都正色道：“这是那里的 话？}
}
{
    \eR{[English source not present in the checked-in Bencraft/Joly files.]}
}
{
}

\Verse{173}
{
    \zR{玩是玩，笑是笑，这个事非同儿戏，你可别混说。}
}
{
    \eR{[English source not present in the checked-in Bencraft/Joly files.]}
}
{
}

\Verse{174}
{
    \zR{你自己昏了心了，想想罢， 想想搁在那里了？}
}
{
    \eR{[English source not present in the checked-in Bencraft/Joly files.]}
}
{
}

\Verse{175}
{
    \zR{这会子又混赖人了！”}
}
{
    \eR{[English source not present in the checked-in Bencraft/Joly files.]}
}
{
}

\Verse{176}
{
    \zR{袭人见他这般光景不像是玩话，便着急道： “皇天菩萨!小祖宗!你到底撂在那里了？”}
}
{
    \eR{[English source not present in the checked-in Bencraft/Joly files.]}
}
{
}

\Verse{177}
{
    \zR{宝玉道：“我记的明明儿放在炕桌上， 你们到底找啊。”}
}
{
    \eR{[English source not present in the checked-in Bencraft/Joly files.]}
}
{
}

\Verse{178}
{
    \zR{袭人麝月等也不敢叫人知道，大家偷偷儿的各处搜寻。}
}
{
    \eR{[English source not present in the checked-in Bencraft/Joly files.]}
}
{
}

\Verse{179}
{
    \zR{闹了大半天，毫无影响， 甚至翻箱倒笼，实在没处去找，便疑到方才这些人进来，不知谁检了去了。}
}
{
    \eR{[English source not present in the checked-in Bencraft/Joly files.]}
}
{
}

\Verse{180}
{
    \zR{袭人说 道：“进来的，谁不知道这玉是性命似的东西呢？}
}
{
    \eR{[English source not present in the checked-in Bencraft/Joly files.]}
}
{
}

\Verse{181}
{
    \zR{谁敢检了去!你们好歹先别声张， 快到各处问去。}
}
{
    \eR{[English source not present in the checked-in Bencraft/Joly files.]}
}
{
}

\Verse{182}
{
    \zR{若有姐妹们检着和我们玩呢，你们给他磕个头，要了来；要是小丫 头们偷了去，问出来，也不回上头，不论做些什么送他换了来，都使得的。}
}
{
    \eR{[English source not present in the checked-in Bencraft/Joly files.]}
}
{
}

\Verse{183}
{
    \zR{这可不 是小事，真要丢了这个，比丢了宝二爷的还利害呢！”}
}
{
    \eR{[English source not present in the checked-in Bencraft/Joly files.]}
}
{
}

\Verse{184}
{
    \zR{麝月秋纹刚要往外走，袭人 又赶出来嘱咐道：“头里在这里吃饭的倒别先问去。}
}
{
    \eR{[English source not present in the checked-in Bencraft/Joly files.]}
}
{
}

\Verse{185}
{
    \zR{找不成，再惹出些风波来，更 不好了。”}
}
{
    \eR{[English source not present in the checked-in Bencraft/Joly files.]}
}
{
}

\Verse{186}
{
    \zR{麝月等依言，分头各处追问。}
}
{
    \eR{[English source not present in the checked-in Bencraft/Joly files.]}
}
{
}

\Verse{187}
{
    \zR{人人不晓，个个惊疑。}
}
{
    \eR{[English source not present in the checked-in Bencraft/Joly files.]}
}
{
}

\Verse{188}
{
    \zR{二人连忙回来，俱 目瞪口呆，面面相窥。}
}
{
    \eR{[English source not present in the checked-in Bencraft/Joly files.]}
}
{
}

\Verse{189}
{
    \zR{宝玉也吓怔了，袭人急的只是干哭。}
}
{
    \eR{[English source not present in the checked-in Bencraft/Joly files.]}
}
{
}

\Verse{190}
{
    \zR{找是没处找，回又不敢 回，怡红院里的人吓的一个个像木雕泥塑一般。}
}
{
    \eR{[English source not present in the checked-in Bencraft/Joly files.]}
}
{
}

\Verse{191}
{
    \zR{大家正在发呆，只见各处知道的都来了。}
}
{
    \eR{[English source not present in the checked-in Bencraft/Joly files.]}
}
{
}

\Verse{192}
{
    \zR{探春叫把园门关上，先叫个老婆子带 着两个丫头，再往各处去寻去；一面又叫告诉众人：“若谁找出来，重重的赏他。”}
}
{
    \eR{[English source not present in the checked-in Bencraft/Joly files.]}
}
{
}

\Verse{193}
{
    \zR{大家头宗要脱干系，二宗听见重赏，不顾命的混找了一遍，甚至于茅厕里都找到了。}
}
{
    \eR{[English source not present in the checked-in Bencraft/Joly files.]}
}
{
}

\Verse{194}
{
    \zR{谁知那块玉竟像绣花针儿一般，找了一天，总无影响。}
}
{
    \eR{[English source not present in the checked-in Bencraft/Joly files.]}
}
{
}

\Verse{195}
{
    \zR{李纨急了，说：“这件事不 是玩的，我要说句无礼的话了。”}
}
{
    \eR{[English source not present in the checked-in Bencraft/Joly files.]}
}
{
}

\Verse{196}
{
    \zR{众人道：“什么话？”}
}
{
    \eR{[English source not present in the checked-in Bencraft/Joly files.]}
}
{
}

\Verse{197}
{
    \zR{李纨道：“事情到了这里 也顾不得了。}
}
{
    \eR{[English source not present in the checked-in Bencraft/Joly files.]}
}
{
}

\Verse{198}
{
    \zR{现在园里除了宝玉，都是女人。}
}
{
    \eR{[English source not present in the checked-in Bencraft/Joly files.]}
}
{
}

\Verse{199}
{
    \zR{要求各位姐姐、妹妹、姑娘都要叫跟 来的丫头脱了衣服，大家搜一搜。}
}
{
    \eR{[English source not present in the checked-in Bencraft/Joly files.]}
}
{
}

\Verse{200}
{
    \zR{若没有，再叫丫头们去搜那些老婆子并粗使的丫 头，不知使得使不得？”}
}
{
    \eR{[English source not present in the checked-in Bencraft/Joly files.]}
}
{
}

\Verse{201}
{
    \zR{大家说道：“这话也说的有理。}
}
{
    \eR{[English source not present in the checked-in Bencraft/Joly files.]}
}
{
}

\Verse{202}
{
    \zR{现在人多手乱，鱼龙混杂， 倒是这么着，他们也洗洗清。”}
}
{
    \eR{[English source not present in the checked-in Bencraft/Joly files.]}
}
{
}

\Verse{203}
{
    \zR{探春独不言语。}
}
{
    \eR{[English source not present in the checked-in Bencraft/Joly files.]}
}
{
}

\Verse{204}
{
    \zR{那些丫头们也都愿意洗净自己。}
}
{
    \eR{[English source not present in the checked-in Bencraft/Joly files.]}
}
{
}

\Verse{205}
{
    \zR{先 是平儿起，平儿说道：“打我先搜起。”}
}
{
    \eR{[English source not present in the checked-in Bencraft/Joly files.]}
}
{
}

\Verse{206}
{
    \zR{于是各人自己解怀。}
}
{
    \eR{[English source not present in the checked-in Bencraft/Joly files.]}
}
{
}

\Verse{207}
{
    \zR{李纨一气儿混搜。}
}
{
    \eR{[English source not present in the checked-in Bencraft/Joly files.]}
}
{
}

\Verse{208}
{
    \zR{探 春嗔着李纨道：“大嫂子，你也学那起不成材料的样子来了!那个人既偷了去还肯 藏在身上？}
}
{
    \eR{[English source not present in the checked-in Bencraft/Joly files.]}
}
{
}

\Verse{209}
{
    \zR{况且这件东西，在家里是宝，到了外头不知道的是废物，偷他做什么？}
}
{
    \eR{[English source not present in the checked-in Bencraft/Joly files.]}
}
{
}

\Verse{210}
{
    \zR{我 想来必是有人使促狭。”}
}
{
    \eR{[English source not present in the checked-in Bencraft/Joly files.]}
}
{
}

\Verse{211}
{
    \zR{众人听说，又见环儿不在这里，昨儿是他满屋里乱跑，都疑到他身上，只是不 肯说出来。}
}
{
    \eR{[English source not present in the checked-in Bencraft/Joly files.]}
}
{
}

\Verse{212}
{
    \zR{探春又道：“使促狭的只有环儿。}
}
{
    \eR{[English source not present in the checked-in Bencraft/Joly files.]}
}
{
}

\Verse{213}
{
    \zR{你们叫个人去悄悄的叫了他来，背地 里哄着他，叫他拿出来，然后吓着他叫他别声张就完了。”}
}
{
    \eR{[English source not present in the checked-in Bencraft/Joly files.]}
}
{
}

\Verse{214}
{
    \zR{大家点头。}
}
{
    \eR{[English source not present in the checked-in Bencraft/Joly files.]}
}
{
}

\Verse{215}
{
    \zR{李纨便向平 儿道：“这件事还得你去才弄的明白。”}
}
{
    \eR{[English source not present in the checked-in Bencraft/Joly files.]}
}
{
}

\Verse{216}
{
    \zR{平儿答应，就赶着去了。}
}
{
    \eR{[English source not present in the checked-in Bencraft/Joly files.]}
}
{
}

\Verse{217}
{
    \zR{不多时，同着贾 环来了。}
}
{
    \eR{[English source not present in the checked-in Bencraft/Joly files.]}
}
{
}

\Verse{218}
{
    \zR{众人假意装出没事的样子，叫人沏了茶，搁在里间屋里。}
}
{
    \eR{[English source not present in the checked-in Bencraft/Joly files.]}
}
{
}

\Verse{219}
{
    \zR{众人故意搭讪走 开，原叫平儿哄他。}
}
{
    \eR{[English source not present in the checked-in Bencraft/Joly files.]}
}
{
}

\Verse{220}
{
    \zR{平儿便笑着向贾环道：“你二哥哥的玉丢了，你瞧见了没有？”}
}
{
    \eR{[English source not present in the checked-in Bencraft/Joly files.]}
}
{
}

\Verse{221}
{
    \zR{贾环便急的紫涨了脸，瞪着眼，说道：“人家丢了东西，你怎么又叫我来查问疑我! 我是犯过案的贼么？”}
}
{
    \eR{[English source not present in the checked-in Bencraft/Joly files.]}
}
{
}

\Verse{222}
{
    \zR{平儿见这样子，倒不敢再问，便又陪笑道：“不是这么说。}
}
{
    \eR{[English source not present in the checked-in Bencraft/Joly files.]}
}
{
}

\Verse{223}
{
    \zR{怕三爷要拿了去吓他们，所以白问问瞧见了没有，好叫他们找。”}
}
{
    \eR{[English source not present in the checked-in Bencraft/Joly files.]}
}
{
}
\ChapterImage{img/chapters/190第94回 失寶玉通靈知奇禍.jpg}


\Verse{224}
{
    \zR{贾环道：“他的 玉在他身上，看见没看见该问他，怎么问我呢？}
}
{
    \eR{[English source not present in the checked-in Bencraft/Joly files.]}
}
{
}

\Verse{225}
{
    \zR{你们都捧着他，得了什么不问我， 丢了东西就来问我！”}
}
{
    \eR{[English source not present in the checked-in Bencraft/Joly files.]}
}
{
}

\Verse{226}
{
    \zR{说着，起身就走。}
}
{
    \eR{[English source not present in the checked-in Bencraft/Joly files.]}
}
{
}

\Verse{227}
{
    \zR{众人不好拦他。}
}
{
    \eR{[English source not present in the checked-in Bencraft/Joly files.]}
}
{
}

\Verse{228}
{
    \zR{这里宝玉倒急了，说道： “都是这劳什子闹事!我也不要他了，你们也不用闹了。}
}
{
    \eR{[English source not present in the checked-in Bencraft/Joly files.]}
}
{
}

\Verse{229}
{
    \zR{环儿一去，必是嚷的满院 里都知道了，这可不是闹事了么？”}
}
{
    \eR{[English source not present in the checked-in Bencraft/Joly files.]}
}
{
}

\Verse{230}
{
    \zR{袭人等急的又哭道：“小祖宗儿，你看这玉丢 了没要紧；要是上头知道了，我们这些人就要粉身碎骨了。”}
}
{
    \eR{[English source not present in the checked-in Bencraft/Joly files.]}
}
{
}

\Verse{231}
{
    \zR{说着，便嚎啕大哭起 来。}
}
{
    \eR{[English source not present in the checked-in Bencraft/Joly files.]}
}
{
}

\Verse{232}
{
    \zR{众人更加着急，明知此事掩饰不来，只得要商议定了话，回来好回贾母诸人。}
}
{
    \eR{[English source not present in the checked-in Bencraft/Joly files.]}
}
{
}

\Verse{233}
{
    \zR{宝玉道：“你们竟也不用商量，硬说我砸了就完了。”}
}
{
    \eR{[English source not present in the checked-in Bencraft/Joly files.]}
}
{
}

\Verse{234}
{
    \zR{平儿道：“我的爷，好轻巧 话儿!上头要问为什么砸的呢？}
}
{
    \eR{[English source not present in the checked-in Bencraft/Joly files.]}
}
{
}

\Verse{235}
{
    \zR{他们也是个死啊。}
}
{
    \eR{[English source not present in the checked-in Bencraft/Joly files.]}
}
{
}

\Verse{236}
{
    \zR{倘或要起砸破的碴儿来，那又怎么 样呢？”}
}
{
    \eR{[English source not present in the checked-in Bencraft/Joly files.]}
}
{
}

\Verse{237}
{
    \zR{宝玉道：“不然，就说我出门丢了。”}
}
{
    \eR{[English source not present in the checked-in Bencraft/Joly files.]}
}
{
}

\Verse{238}
{
    \zR{众人一想：“这句话倒还混的过去， 但只这两天又没上学，又没往别处去。”}
}
{
    \eR{[English source not present in the checked-in Bencraft/Joly files.]}
}
{
}

\Verse{239}
{
    \zR{宝玉道：“怎么没有？}
}
{
    \eR{[English source not present in the checked-in Bencraft/Joly files.]}
}
{
}

\Verse{240}
{
    \zR{大前儿还到临安伯 府里听戏去了呢。}
}
{
    \eR{[English source not present in the checked-in Bencraft/Joly files.]}
}
{
}

\Verse{241}
{
    \zR{就说那日丢的就完了。”}
}
{
    \eR{[English source not present in the checked-in Bencraft/Joly files.]}
}
{
}

\Verse{242}
{
    \zR{探春道：“那也不妥。}
}
{
    \eR{[English source not present in the checked-in Bencraft/Joly files.]}
}
{
}

\Verse{243}
{
    \zR{既是前儿丢的， 为什么当日不来回？”}
}
{
    \eR{[English source not present in the checked-in Bencraft/Joly files.]}
}
{
}

\Verse{244}
{
    \zR{众人正在胡思乱想要装点撒谎，只听见赵姨娘的声儿哭着喊 着走来，说：“你们丢了东西，自己不找，怎么叫人背地里拷问环儿!我把环儿带
        了来，索性交给你们这一起？}
}
{
    \eR{[English source not present in the checked-in Bencraft/Joly files.]}
}
{
}

\Verse{245}
{
    \zR{上水的，该杀该剐随你们罢！”}
}
{
    \eR{[English source not present in the checked-in Bencraft/Joly files.]}
}
{
}

\Verse{246}
{
    \zR{说着将环儿一推，说： “你是个贼，快快的招罢！”}
}
{
    \eR{[English source not present in the checked-in Bencraft/Joly files.]}
}
{
}

\Verse{247}
{
    \zR{气的环儿也哭喊起来。}
}
{
    \eR{[English source not present in the checked-in Bencraft/Joly files.]}
}
{
}

\Verse{248}
{
    \zR{李纨正要劝解，丫头来说：“太太来了。”}
}
{
    \eR{[English source not present in the checked-in Bencraft/Joly files.]}
}
{
}

\Verse{249}
{
    \zR{袭人等此时无地可容。}
}
{
    \eR{[English source not present in the checked-in Bencraft/Joly files.]}
}
{
}

\Verse{250}
{
    \zR{宝玉等赶忙 出来迎接。}
}
{
    \eR{[English source not present in the checked-in Bencraft/Joly files.]}
}
{
}

\Verse{251}
{
    \zR{赵姨娘暂且也不敢作声，跟了出来。}
}
{
    \eR{[English source not present in the checked-in Bencraft/Joly files.]}
}
{
}

\Verse{252}
{
    \zR{王夫人见众人都有惊惶之色，才信 方才听见的话，便道：“那块玉真丢了么？”}
}
{
    \eR{[English source not present in the checked-in Bencraft/Joly files.]}
}
{
}

\Verse{253}
{
    \zR{众人都不敢作声。}
}
{
    \eR{[English source not present in the checked-in Bencraft/Joly files.]}
}
{
}

\Verse{254}
{
    \zR{王夫人走进屋里坐 下，便叫袭人，慌的袭人连忙跪下，含泪要禀。}
}
{
    \eR{[English source not present in the checked-in Bencraft/Joly files.]}
}
{
}

\Verse{255}
{
    \zR{王夫人道：“你起来，快快叫人细 细的找去，一忙乱倒不好了。”}
}
{
    \eR{[English source not present in the checked-in Bencraft/Joly files.]}
}
{
}

\Verse{256}
{
    \zR{袭人哽咽难言。}
}
{
    \eR{[English source not present in the checked-in Bencraft/Joly files.]}
}
{
}

\Verse{257}
{
    \zR{宝玉恐袭人直告诉出来，便说道： “太太，这事不与袭人相干，是我前日到临安伯府里听戏在路上丢了。”}
}
{
    \eR{[English source not present in the checked-in Bencraft/Joly files.]}
}
{
}

\Verse{258}
{
    \zR{王夫人道： “为什么那日不找呢？”}
}
{
    \eR{[English source not present in the checked-in Bencraft/Joly files.]}
}
{
}

\Verse{259}
{
    \zR{宝玉道：“我怕他们知道，没有告诉他们。}
}
{
    \eR{[English source not present in the checked-in Bencraft/Joly files.]}
}
{
}

\Verse{260}
{
    \zR{我叫焙茗等在 外头各处找过的。”}
}
{
    \eR{[English source not present in the checked-in Bencraft/Joly files.]}
}
{
}

\Verse{261}
{
    \zR{王夫人道：“胡说，如今脱换衣服，不是袭人他们伏侍的么？}
}
{
    \eR{[English source not present in the checked-in Bencraft/Joly files.]}
}
{
}

\Verse{262}
{
    \zR{大凡哥儿出门回来，手巾荷包短了，还要个明白，何况这块玉不见了，难道不问 么？”}
}
{
    \eR{[English source not present in the checked-in Bencraft/Joly files.]}
}
{
}

\Verse{263}
{
    \zR{宝玉无言可答。}
}
{
    \eR{[English source not present in the checked-in Bencraft/Joly files.]}
}
{
}

\Verse{264}
{
    \zR{赵姨娘听见，便得意了，忙接口道：“外头丢了东西，也赖 环儿――”话未说完，被王夫人喝道：“这里说这个，你且说那些没要紧的话！”}
}
{
    \eR{[English source not present in the checked-in Bencraft/Joly files.]}
}
{
}

\Verse{265}
{
    \zR{赵姨娘便也不敢言语了。}
}
{
    \eR{[English source not present in the checked-in Bencraft/Joly files.]}
}
{
}

\Verse{266}
{
    \zR{还是李纨探春从实的告诉了王夫人一遍。}
}
{
    \eR{[English source not present in the checked-in Bencraft/Joly files.]}
}
{
}

\Verse{267}
{
    \zR{王夫人也急的眼 中落泪，索性要回明了贾母，去问邢夫人那边来的这些人去。}
}
{
    \eR{[English source not present in the checked-in Bencraft/Joly files.]}
}
{
}

\Verse{268}
{
    \zR{凤姐病中也听见宝玉失玉，知道王夫人过来，料躲不住，便扶了丰儿来到园里。}
}
{
    \eR{[English source not present in the checked-in Bencraft/Joly files.]}
}
{
}

\Verse{269}
{
    \zR{正值王夫人起身要走，凤姐娇怯怯的说：“请太太安。”}
}
{
    \eR{[English source not present in the checked-in Bencraft/Joly files.]}
}
{
}

\Verse{270}
{
    \zR{宝玉等过来问了凤姐好。}
}
{
    \eR{[English source not present in the checked-in Bencraft/Joly files.]}
}
{
}

\Verse{271}
{
    \zR{王夫人因说道：“你也听见了么？}
}
{
    \eR{[English source not present in the checked-in Bencraft/Joly files.]}
}
{
}

\Verse{272}
{
    \zR{这可不是奇事吗？}
}
{
    \eR{[English source not present in the checked-in Bencraft/Joly files.]}
}
{
}

\Verse{273}
{
    \zR{刚才眼错不见就丢了，再找不着。}
}
{
    \eR{[English source not present in the checked-in Bencraft/Joly files.]}
}
{
}

\Verse{274}
{
    \zR{你去想想：打老太太那边的丫头起，至你们平儿，谁的手不稳，谁的心促狭，我要 回了老太太，认真的查出来才好。}
}
{
    \eR{[English source not present in the checked-in Bencraft/Joly files.]}
}
{
}

\Verse{275}
{
    \zR{不然，是断了宝玉的命根子了！”}
}
{
    \eR{[English source not present in the checked-in Bencraft/Joly files.]}
}
{
}

\Verse{276}
{
    \zR{凤姐回道：“咱 们家人多手杂，自古说的，‘知人知面不知心’，那里保的住谁是好的？}
}
{
    \eR{[English source not present in the checked-in Bencraft/Joly files.]}
}
{
}

\Verse{277}
{
    \zR{但只一吵 嚷，已经都知道了，偷玉的人要叫太太查出来，明知是死无葬身之地，他着了急， 反要毁坏了灭口，那时可怎么处呢。}
}
{
    \eR{[English source not present in the checked-in Bencraft/Joly files.]}
}
{
}

\Verse{278}
{
    \zR{据我的糊涂想头，只说宝玉本不爱他，撂丢了， 也没有什么要紧，只要大家严密些，别叫老太太老爷知道。}
}
{
    \eR{[English source not present in the checked-in Bencraft/Joly files.]}
}
{
}

\Verse{279}
{
    \zR{这么说了，暗暗的派人 去各处察访，哄骗出来，那时玉也可得，罪名也可定：不知太太心里怎么样？”}
}
{
    \eR{[English source not present in the checked-in Bencraft/Joly files.]}
}
{
}

\Verse{280}
{
    \zR{王 夫人迟了半日，才说道：“你这话虽也有理，但只是老爷跟前怎么瞒的过呢？”}
}
{
    \eR{[English source not present in the checked-in Bencraft/Joly files.]}
}
{
}

\Verse{281}
{
    \zR{便 叫环儿来说道：“你二哥哥的玉丢了，白问了你一句，怎么你就乱嚷？}
}
{
    \eR{[English source not present in the checked-in Bencraft/Joly files.]}
}
{
}

\Verse{282}
{
    \zR{要是嚷破了， 人家把那个毁坏了，我看你活得活不得！”}
}
{
    \eR{[English source not present in the checked-in Bencraft/Joly files.]}
}
{
}

\Verse{283}
{
    \zR{贾环吓得哭道：“我再不敢嚷了。”}
}
{
    \eR{[English source not present in the checked-in Bencraft/Joly files.]}
}
{
}

\Verse{284}
{
    \zR{赵 姨娘听了，那里还敢言语。}
}
{
    \eR{[English source not present in the checked-in Bencraft/Joly files.]}
}
{
}

\Verse{285}
{
    \zR{王夫人便吩咐众人道：“想来自然有没找到的地方儿。}
}
{
    \eR{[English source not present in the checked-in Bencraft/Joly files.]}
}
{
}

\Verse{286}
{
    \zR{好端端的在家里的，还怕他飞到那里去不成？}
}
{
    \eR{[English source not present in the checked-in Bencraft/Joly files.]}
}
{
}

\Verse{287}
{
    \zR{只是不许声张。}
}
{
    \eR{[English source not present in the checked-in Bencraft/Joly files.]}
}
{
}

\Verse{288}
{
    \zR{限袭人三天内给我找 出来。}
}
{
    \eR{[English source not present in the checked-in Bencraft/Joly files.]}
}
{
}

\Verse{289}
{
    \zR{要是三天找不着，只怕也瞒不住，大家那就不用过安静日子了！”}
}
{
    \eR{[English source not present in the checked-in Bencraft/Joly files.]}
}
{
}

\Verse{290}
{
    \zR{说着，便 叫凤姐儿跟到邢夫人那边，商议踩缉不提。}
}
{
    \eR{[English source not present in the checked-in Bencraft/Joly files.]}
}
{
}

\Verse{291}
{
    \zR{这里李纨等纷纷议论，便传唤看园子的一干人来，叫把园门锁上，快传林之孝 家的来，悄悄儿的告诉了他，叫他：“吩咐前后门上：三天之内，不论男女下人，
        从里头可以走动，要出去时，一概不许放出。}
}
{
    \eR{[English source not present in the checked-in Bencraft/Joly files.]}
}
{
}

\Verse{292}
{
    \zR{只说里头丢了东西，等这件东西有了 着落，然后放人出来。”}
}
{
    \eR{[English source not present in the checked-in Bencraft/Joly files.]}
}
{
}

\Verse{293}
{
    \zR{林之孝家的答应了“是”，因说：“前儿奴才家里也丢了 一件不要紧的东西，林之孝必要明白，上街去找了一个测字的。}
}
{
    \eR{[English source not present in the checked-in Bencraft/Joly files.]}
}
{
}

\Verse{294}
{
    \zR{那人叫做什么刘铁 嘴，测了一个字，说的很明白，回来按着一找，就找着了。”}
}
{
    \eR{[English source not present in the checked-in Bencraft/Joly files.]}
}
{
}

\Verse{295}
{
    \zR{袭人听见，便央及林 家的道：“好林奶奶，出去快求林大爷替我们问问去。”}
}
{
    \eR{[English source not present in the checked-in Bencraft/Joly files.]}
}
{
}

\Verse{296}
{
    \zR{那林之孝家的答应着出去 了。}
}
{
    \eR{[English source not present in the checked-in Bencraft/Joly files.]}
}
{
}

\Verse{297}
{
    \zR{邢岫烟道：“若说那外头测字打卦的，是不中用的。}
}
{
    \eR{[English source not present in the checked-in Bencraft/Joly files.]}
}
{
}

\Verse{298}
{
    \zR{我在南边闻妙玉能扶乩， 何不烦他问一问？}
}
{
    \eR{[English source not present in the checked-in Bencraft/Joly files.]}
}
{
}

\Verse{299}
{
    \zR{况且我听见说，这块玉原有仙机，想来问的出来。”}
}
{
    \eR{[English source not present in the checked-in Bencraft/Joly files.]}
}
{
}

\Verse{300}
{
    \zR{众人都诧异 道：“咱们常见的，从没有听他说起。”}
}
{
    \eR{[English source not present in the checked-in Bencraft/Joly files.]}
}
{
}

\Verse{301}
{
    \zR{麝月便忙问岫烟道：“想来别人求他是不 肯的，好姑娘，我给姑娘磕个头，求姑娘就去!若问出来了，我一辈子总不忘你的 恩。”}
}
{
    \eR{[English source not present in the checked-in Bencraft/Joly files.]}
}
{
}

\Verse{302}
{
    \zR{说着，赶忙就要磕下头去，岫烟连忙拦住。}
}
{
    \eR{[English source not present in the checked-in Bencraft/Joly files.]}
}
{
}

\Verse{303}
{
    \zR{黛玉等也都怂恿着岫烟速往栊翠 庵去。}
}
{
    \eR{[English source not present in the checked-in Bencraft/Joly files.]}
}
{
}

\Verse{304}
{
    \zR{一面林之孝家的进来说道：“姑娘们大喜!林之孝测了字回来，说这玉是丢不 了的，将来横竖有人送还来的。”}
}
{
    \eR{[English source not present in the checked-in Bencraft/Joly files.]}
}
{
}

\Verse{305}
{
    \zR{众人听了，也都半信半疑，惟有袭人麝月喜欢的 了不得。}
}
{
    \eR{[English source not present in the checked-in Bencraft/Joly files.]}
}
{
}

\Verse{306}
{
    \zR{探春便问：“测的是什么字？”}
}
{
    \eR{[English source not present in the checked-in Bencraft/Joly files.]}
}
{
}

\Verse{307}
{
    \zR{林之孝家的道：“他的话多，奴才也学不 上来。}
}
{
    \eR{[English source not present in the checked-in Bencraft/Joly files.]}
}
{
}

\Verse{308}
{
    \zR{记得是拈了个赏人东西的‘赏’字。}
}
{
    \eR{[English source not present in the checked-in Bencraft/Joly files.]}
}
{
}

\Verse{309}
{
    \zR{那刘铁嘴也不问，便说：‘丢了东西不 是？’”}
}
{
    \eR{[English source not present in the checked-in Bencraft/Joly files.]}
}
{
}

\Verse{310}
{
    \zR{李纨道：“这就算好。”}
}
{
    \eR{[English source not present in the checked-in Bencraft/Joly files.]}
}
{
}

\Verse{311}
{
    \zR{林之孝家的道：“他还说：‘“赏”字上头一个 “小”字，底下一个“口”字，这件东西，很可嘴里放得，必是个珠子宝石。’”}
}
{
    \eR{[English source not present in the checked-in Bencraft/Joly files.]}
}
{
}

\Verse{312}
{
    \zR{众人听了，夸赞道：“真是神仙!往下怎么说？”}
}
{
    \eR{[English source not present in the checked-in Bencraft/Joly files.]}
}
{
}

\Verse{313}
{
    \zR{林之孝家的道：“他说：‘底下 “贝”字拆开，不成一个“见”字，可不是“不见”了？’}
}
{
    \eR{[English source not present in the checked-in Bencraft/Joly files.]}
}
{
}

\Verse{314}
{
    \zR{因上头拆了‘当’字， 叫快到当铺里找去。}
}
{
    \eR{[English source not present in the checked-in Bencraft/Joly files.]}
}
{
}

\Verse{315}
{
    \zR{‘赏’字加一‘人’字，可不是‘？’}
}
{
    \eR{[English source not present in the checked-in Bencraft/Joly files.]}
}
{
}

\Verse{316}
{
    \zR{字？}
}
{
    \eR{[English source not present in the checked-in Bencraft/Joly files.]}
}
{
}

\Verse{317}
{
    \zR{只要找着当铺就有 人，有了人便赎了来，可不是偿还了吗？”}
}
{
    \eR{[English source not present in the checked-in Bencraft/Joly files.]}
}
{
}

\Verse{318}
{
    \zR{众人道：“既这么着，就先往左近找起。}
}
{
    \eR{[English source not present in the checked-in Bencraft/Joly files.]}
}
{
}

\Verse{319}
{
    \zR{横竖几个当铺都找遍了，少不得就有了。}
}
{
    \eR{[English source not present in the checked-in Bencraft/Joly files.]}
}
{
}

\Verse{320}
{
    \zR{咱们有了东西，再问人就容易了。”}
}
{
    \eR{[English source not present in the checked-in Bencraft/Joly files.]}
}
{
}

\Verse{321}
{
    \zR{李纨 道：“只要东西，那怕不问人都使得。}
}
{
    \eR{[English source not present in the checked-in Bencraft/Joly files.]}
}
{
}

\Verse{322}
{
    \zR{林嫂子你去，就把测字的话快告诉了二奶奶， 回了太太，也叫太太放心。}
}
{
    \eR{[English source not present in the checked-in Bencraft/Joly files.]}
}
{
}

\Verse{323}
{
    \zR{就叫二奶奶快派人查去。”}
}
{
    \eR{[English source not present in the checked-in Bencraft/Joly files.]}
}
{
}

\Verse{324}
{
    \zR{林家的答应了便走。}
}
{
    \eR{[English source not present in the checked-in Bencraft/Joly files.]}
}
{
}

\Verse{325}
{
    \zR{众人略安了一点儿神，呆呆的等岫烟回来。}
}
{
    \eR{[English source not present in the checked-in Bencraft/Joly files.]}
}
{
}

\Verse{326}
{
    \zR{正呆等时，只见跟宝玉的焙茗在门 外招手儿，叫小丫头子快出来。}
}
{
    \eR{[English source not present in the checked-in Bencraft/Joly files.]}
}
{
}

\Verse{327}
{
    \zR{那小丫头赶忙的出去了。}
}
{
    \eR{[English source not present in the checked-in Bencraft/Joly files.]}
}
{
}

\Verse{328}
{
    \zR{焙茗便说道：“你快进去 告诉我们二爷和里头太太、奶奶、姑娘们，天大的喜事！”}
}
{
    \eR{[English source not present in the checked-in Bencraft/Joly files.]}
}
{
}

\Verse{329}
{
    \zR{那小丫头子道：“你快 说罢，怎么这么累赘？”}
}
{
    \eR{[English source not present in the checked-in Bencraft/Joly files.]}
}
{
}

\Verse{330}
{
    \zR{焙茗笑着拍手道：“我告诉姑娘，姑娘进去回了，咱们两 个人都得赏钱呢。}
}
{
    \eR{[English source not present in the checked-in Bencraft/Joly files.]}
}
{
}

\Verse{331}
{
    \zR{你打量是什么事情？}
}
{
    \eR{[English source not present in the checked-in Bencraft/Joly files.]}
}
{
}

\Verse{332}
{
    \zR{宝二爷的那块玉呀，我得了准信儿来了。”}
}
{
    \eR{[English source not present in the checked-in Bencraft/Joly files.]}
}
{
}

\Verse{333}
{
    \zR{未知如何，下回分解。}
}
{
    \eR{[English source not present in the checked-in Bencraft/Joly files.]}
}
{
}

\Verse{334}
{
    \zR{(本章完)}
}
{
    \eR{[English source not present in the checked-in Bencraft/Joly files.]}
}
{
}
